\section*{Chapter \ref{ch:rs:overfitting} --- Overfitting}

\begin{solution}{exo:rs:overfitting:1}
$\omega = 380/1\,001 = 0.380$, $\lambda = \ln(0.380/0.620) = -0.49$: below the median, so the split counts toward the PBO.
\end{solution}

\begin{solution}{exo:rs:overfitting:2}
$\binom{10}{2}\times 2/10 = 45 \times 2/10 = 9$ paths; $\binom{6}{3} \times 3/6 = 10$.
\end{solution}

\begin{solution}{exo:rs:overfitting:3}
$2\ln 100/1^2 = 9.2$ years.
\end{solution}

\begin{solution}{exo:rs:overfitting:4}
A label observed on day $d$ spans days $d$ to $d + 9$. The test labels span days 200 to 308. \omterm{def:rs:overfitting:purging}{Purging} removes the training days whose labels touch that span:
191 to 199 (their labels reach into day 200 and after) and 300 to 308 (they start inside it). An \omterm{def:rs:overfitting:purging}{embargo} of 5 days also removes days 309 to 313.
\end{solution}

\begin{solution}{exo:rs:overfitting:5}
Correlated systems share most of their noise; the \omterm{def:rs:overfitting:overfitting}{in-sample} winner is the one whose idiosyncratic part fitted the \omterm{def:rs:overfitting:overfitting}{in-sample} half best. With the halves
drawn symmetrically, the same idiosyncratic part tends to fit the other half worse than the common part does for the others, so the winner falls below
the median more often than not: selection anticorrelates with \omterm{def:rs:overfitting:overfitting}{out-of-sample} rank (the slopes of $-0.89$ and $-0.91$). With independent strategies and no
skill the PBO is near one half.
\end{solution}

\begin{solution}{exo:rs:overfitting:6}
The trend is common to all systems: it lifts them all, by about the same amount, and does not distinguish among them. What distinguishes them in sample is
noise, which does not persist; the winner of a search over noise on top of a common trend earns the trend minus its selection penalty (0.14 against an
average of 0.26).
\end{solution}

\begin{solution}{exo:rs:overfitting:7}
\texttt{rs\_overfit.pbo\_hundred}: 0.69 for the first 100 trend systems, 0.52 for 100 independent noise strategies. Independent strategies with no skill give a
PBO near one half, as selection among pure noise should; correlated strategies push it above, because the winner's edge is the part of its noise the others
do not share.
\end{solution}

\begin{solution}{exo:rs:overfitting:8}
Shuffled folds put overlapping labels and serially correlated features on both sides of every split: the model is scored on days whose labels it has, in
effect, already seen. On the chapter's noise, shuffled folds gave an $R^2$ of 0.90 and every time-respecting scheme a negative one. Use contiguous folds,
purge overlapping labels, \omterm{def:rs:overfitting:purging}{embargo} after the test folds, and tune on folds that precede the test period.
\end{solution}

\begin{solution}{pb:rs:overfitting:1}
\begin{enumerate}
\item Noise: 0.44 in sample, $-0.32$ out of sample. Trend: 0.77 and 0.14.
\item Noise: median in sample 0.17, average out of sample 0.20. Trend: 0.50 and 0.26.
\item 0.77; the systems are highly correlated (their Sharpe ratios spread by 0.076), so the search amounts to few independent trials, and the expected best
given that spread is 0.25 above the common level.
\item The trend lifts every system (median 0.50). The search adds a ranking that carries some information across the systems (correlation 0.33,
the fifty best earn 0.29) but picks its single winner mostly for noise.
\item From block sums and sums of squares: each split is a set of blocks, so its means and variances are matrix products.
\item Noise: PBO 0.82, slope $-0.89$. Trend: 0.74, slope $-0.91$.
\item Selection is worse than choosing at random: the \omterm{def:rs:overfitting:overfitting}{in-sample} winner is more often than not below the median out of sample.
\item Near zero: it wins in sample and stays on top out of sample in almost every split (the build's test finds below 1\%).
\item 42.4, 10.6 and 3.3 years.
\item The standard error of a Sharpe ratio falls with the square root of the history; to halve the target the error must halve, which takes four times the
history.
\item Shuffled 0.90, contiguous $-0.63$, purged $-0.57$, purged and embargoed $-0.93$.
\item Neighbouring days, whose labels share periods, fall into both sets; \omterm{def:rs:overfitting:purging}{purging} removes the training days whose labels overlap the test labels, the \omterm{def:rs:overfitting:purging}{embargo} the
days just after the test set, whose features still carry its information.
\item $-0.04$ on noise, 0.29 with the trend.
\item \textbf{Named result.} The search's PBO is 0.82 on noise and 0.74 with a real trend; the chosen system's \omterm{def:rs:overfitting:overfitting}{out-of-sample} Sharpe ratio is $-0.32$ and 0.14, below
the average system's 0.20 and 0.26 in each world, while annual walk-forward re-selection earns $-0.04$ and 0.29. The deflated expectation of the winner is the
average system's, not its own \omterm{def:rs:overfitting:overfitting}{in-sample} 0.44 or 0.77.
\item The walk-forward procedure (or the average of the plateau), which earned the trend without paying for the choice of a point.
\item The thousand parameter pairs as trials, the data span, the selection rule, the PBO, the \omterm{def:rs:overfitting:minbtl}{minimum backtest length} and the deflated Sharpe ratio, before any
\omterm{def:rs:overfitting:overfitting}{out-of-sample} look.
\item A reason for the trend (hedging pressure, slow diffusion) predicts that the broad class of systems should work, not a particular window, and points at the
plateau rather than the peak.
\item A haircut of the reported Sharpe ratio for the number of tests, a single number to set beside the PBO's probability.
\item Many \omterm{def:rs:overfitting:overfitting}{out-of-sample} paths instead of one, so the procedure's performance is a distribution.
\item Choosing, among many \omterm{def:rs:vectorised-backtests:backtest}{backtests} on limited history, the one whose success is the sample's rather than the strategy's.
\end{enumerate}
\end{solution}

\begin{solution}{iq:rs:overfitting:1}
Record the 500 trials; compute the deflated Sharpe ratio of the best (Book~4, chapter~12) and the PBO of the search by CSCV; check the \omterm{def:rs:overfitting:minbtl}{minimum backtest length} for
500 trials at a Sharpe ratio of 2; look at the neighbourhood of the best (plateau or spike); find a mechanism; then test on data untouched by the search, or trade
small (chapter~21).
\end{solution}

\begin{solution}{iq:rs:overfitting:2}
Observations are serially correlated and labels overlap, so shuffled folds leak test information into training. Use contiguous folds with \omterm{def:rs:overfitting:purging}{purging} of overlapping
labels and an \omterm{def:rs:overfitting:purging}{embargo} after each test fold, walk-forward for the procedure, and \omterm{def:rs:overfitting:purging}{combinatorial purged cross-validation} for a distribution of paths.
\end{solution}

\begin{solution}{iq:rs:overfitting:3}
The PBO is the probability that the \omterm{def:rs:overfitting:overfitting}{in-sample} winner of a search performs below the median out of sample. CSCV cuts the history into blocks, takes every
symmetric split into \omterm{def:rs:overfitting:overfitting}{in-sample} and \omterm{def:rs:overfitting:overfitting}{out-of-sample} halves, picks the winner in sample and ranks it out of sample; the share of splits with a rank below the median
estimates the PBO.
\end{solution}

\begin{solution}{iq:rs:overfitting:4}
About $2\ln 1\,000 = 13.8$ years by the approximation, 10.6 with the exact expected maximum of independent trials; correlated trials need less, and a trustworthy
answer needs more (to make the chance level well below the target, not equal to it).
\end{solution}

\begin{solution}{iq:rs:overfitting:5}
A \omterm{def:rs:the-research-process:log}{research log} with every trial counted; protocols fixed before results; a holdout guarded by someone else; review of searches by their PBO and deflated statistics;
rewarding mechanisms and negative results; limits on the number of variants a project may try per year of data.
\end{solution}

\begin{solution}{iq:rs:overfitting:6}
Null Sharpe ratio estimates over $y$ years are about normal with variance $1/y$; the expected maximum of $N$ of them is about $\sqrt{2\ln N}/\sqrt y$; setting it equal to the
target $s$ gives $y = 2\ln N/s^2$: logarithmic in the number of trials, inverse square in the target.
\end{solution}
